\section{A General Nonviscous Damping Formulation}
Consider the general equation of motion of a dynamic system under an external load $\bp\left(t\right)$,
\begin{gather}
    \br\left(\bu,\bv,\ba\right)=\bp\left(t\right),
\end{gather}
where $\br=\br\left(\bu,\bv,\ba\right)$ is the overall resistance of the system, accounting for contributions from displacement $\bu$, velocity $\bv$ and acceleration $\ba$.
Assuming that the contributions of $\bu$, $\bv=\dot{\bu}$, and $\ba=\ddot{\bu}$ are separable, there is no coupling among them, the resistance $\br$ can be then divided into three components
\begin{gather}
    \br=\br_{\bu}+\br_{\bv}+\br_{\ba},
\end{gather}
where $\br_{\baleph}=\br_{\baleph}\left(\baleph\right)$ is a function of $\baleph\in\left\{\bu,\bv,\ba\right\}$.
Additional energy-dissipating terms can be introduced, for example,
\begin{gather}\label{eq:general_eom}
\boxed{\br+\bbf=\bp,}
\end{gather}
where $\bbf$ can be additively decomposed as
\begin{gather}
    \bbf=\sum\bbf_j.
\end{gather}
Each term $\bbf_j$ is assumed to be governed by a first-order differential equation such that
\begin{gather}\label{eq:general_de}
    g\left(\dot{\bbf}_j,\bbf_j,\by_j\right)=\mb{0},
\end{gather}
where $\by_j=\by_j\left(\bu,\bv,\ba\right)$ is an internal field.

For simplicity, throughout this work we assume that each $\bbf_j$ is governed by the following linear differential equation:
\begin{gather}\label{eq:linear_de}
    \dot{\bbf}_j=-s_j\bbf_j+m_j\by_j
\end{gather}
where $s_j\in\mathbb{R}$ and $m_j\in\mathbb{R}$ are distinct real scalar parameters that control the response.
Both can be extended to the complex domain $\mathbb{C}$ subject to additional conditions, as discussed below.
To ensure stability, the poles $s_j$ must satisfy the necessary condition $\Re\left(s_j\right)>0$.
It is also assumed that the initial condition for $\bbf_j$ is trivial; that is, $\bbf_j\left(0\right)=\mb{0}$.
\subsection{Formulation 1}\label{sec:u_dep}
For ease of discussion, we consider the linear version of \eqsref{eq:general_eom} under free vibration ($\bp=\mb{0}$), which is
\begin{gather}
    \left\{\begin{array}{l}
        \mb{M}\ddot{\bu}+\mb{K}\bu+\sum\bbf_j=\mb{0}, \\
        \dot{\bbf}_j=-s_j\bbf_j+m_j\mb{K}\bu,
    \end{array}\right.
\end{gather}
with $\by_j=\mb{K}\bu$ and real parameters $s_j$ and $m_j$.
Assuming harmonic solutions for both $\bu$ and $\bbf_j$,
\begin{gather}\label{eq:harmonic_solution}
    \bu\left(t\right)=\mb{U}\left(\omega\right)\exp\left(i\omega{}t\right),\qquad
    \bbf_j\left(t\right)=\mb{F}_j\left(\omega\right)\exp\left(i\omega{}t\right),
\end{gather}
the second expression can be written in the frequency domain as
\begin{gather}
    i\omega\mb{F}_j=-s_j\mb{F}_j+m_j\mb{K}\mb{U},
\end{gather}
so that $\mb{F}_j=\dfrac{m_j}{s_j+i\omega}\mb{K}\mb{U}$.
The characteristic equation can then be derived as
\begin{gather}\label{eq:disp_cha}
    \det\left(\left(1+\sum\dfrac{m_j}{s_j+i\omega}\right)\mb{K}-\omega^2\mb{M}\right)=0.
\end{gather}
The dynamic stiffness $\mb{K}^*$ is
\begin{gather}
    \mb{K}^*=\left(1+\sum\dfrac{m_j}{s_j+i\omega}\right)\mb{K}
\end{gather}
with storage stiffness $\mb{K}'$ and loss stiffness $\mb{K}''$
\begin{gather}\label{eq:u_stiffnesses}
    \mb{K}'=\left(1+\sum\dfrac{s_jm_j}{s_j^2+\omega^2}\right)\mb{K},\qquad
    \mb{K}''=-\sum\dfrac{m_j\omega}{s_j^2+\omega^2}\mb{K}
\end{gather}
such that $\mb{K}^*=\mb{K}'+\mb{K}''i$.
%The ratio $\eta_u$ between them is
%\begin{gather}
%\eta_u=-\dfrac{\sum\dfrac{m_j\omega}{s_j^2+\omega^2}}{1+\sum\dfrac{s_jm_j}{s_j^2+\omega^2}}.
%\end{gather}
The limiting cases are
\begin{gather}
    \lim\limits_{\omega\rightarrow0}\mb{K}'=\left(1+\sum\dfrac{m_j}{s_j}\right)\mb{K},\qquad
    \lim\limits_{\omega\rightarrow\infty}\mb{K}'=\mb{K}.
\end{gather}
In this formulation, $m_j<0$ is required to ensure that the system is dissipative.
Physically, $\bbf_j$ represent filtered hysteretic resistance forces and must be exerted against motion in order to dissipate energy.
Accordingly, the storage stiffness $\mb{K}'$ softens at low frequencies (since $m_j<0$ and $s_j>0$) and approaches the original stiffness $\mb{K}$ at high frequencies.

Repeating this procedure with $\by_j=\bu$ gives the following stiffnesses:
\begin{gather}
    \mb{K}'=\mb{K}+\sum\dfrac{s_jm_j}{s_j^2+\omega^2}\mb{I},\qquad
    \mb{K}''=-\sum\dfrac{m_j\omega}{s_j^2+\omega^2}\mb{I}.
\end{gather}
With this choice, the loss stiffness $\mb{K}''$ no longer depends on the static stiffness $\mb{K}$.
\subsection{Complex Realm Extension}
The above analyses assume $s_j$ and $m_j$ are real scalars.
If $m_j$ and $s_j$ come in conjugate pairs, for example,
\begin{gather}
    \boxed{
    m_j\in\mathbb{C},\qquad
    s_j\in\mathbb{C},\qquad
    m_{j+1}=\overline{m_j},\qquad
    s_{j+1}=\overline{s_j},}
\end{gather}
then the corresponding expressions for storage and loss stiffnesses still hold.
For example, the real part is
\begin{multline}
\Re\left(\dfrac{m_j}{s_j+i\omega}+\dfrac{m_{j+1}}{s_{j+1}+i\omega}\right)
=\Re\left(\dfrac{m_j}{s_j+i\omega}+\dfrac{\overline{m_j}}{\overline{s_j}+i\omega}\right)
\\=\Re\left(\dfrac{m_j}{s_j+i\omega}+\overline{\dfrac{m_j}{s_j-i\omega}}\right)
=\Re\left(\dfrac{m_j}{s_j+i\omega}+\dfrac{m_j}{s_j-i\omega}\right)
=\Re\left(\dfrac{2s_jm_j}{s^2_j+\omega^2}\right)
\\=\Re\left(\dfrac{s_jm_j}{s_j^2+\omega^2}+\dfrac{s_jm_j}{s_j^2+\omega^2}\right)
=\Re\left(\dfrac{s_jm_j}{s_j^2+\omega^2}+\overline{\dfrac{s_jm_j}{s_j^2+\omega^2}}\right)
=\Re\left(\dfrac{s_jm_j}{s_j^2+\omega^2}+\dfrac{s_{j+1}m_{j+1}}{s_{j+1}^2+\omega^2}\right).
\end{multline}
This implies that the parameters $s_j$ and $m_j$ can be extended to the complex domain, provided that they come in conjugate pairs.
In the following, we therefore assume $s_j\in\mathbb{C}$ and $m_j\in\mathbb{C}$ with this additional constraint implied.
Again, a stable system requires $\Re\left(s_j\right)>0$.
\subsection{Formulation 2}\label{sec:a_dep}
Alternatively, one can pick $\by_j=\mb{M}\ba$.
The system becomes
\begin{gather}
    \left\{\begin{array}{l}
        \mb{M}\ddot{\bu}+\mb{K}\bu+\sum\bbf_j=\mb{0}, \\
        \dot{\bbf}_j=-s_j\bbf_j+m_j\mb{M}\ba.
    \end{array}\right.
\end{gather}
Adopting a similar approach, the characteristic equation can be derived as
\begin{gather}\label{eq:acc_cha}
    \det\left(\dfrac{1}{1+\sum\dfrac{m_j}{s_j+i\omega}}\mb{K}-\omega^2\mb{M}\right)=0.
\end{gather}
The magnitude of the loss ratio remains the same because the scalar coefficient is the reciprocal of that in the previous case.
In this formulation, a dissipative system requires $\Re\left(m_j\right)>0$.

It is possible to show that \eqsref{eq:acc_cha} is in fact identical to \eqsref{eq:disp_cha}.
Rewriting the scalar coefficient in matrix form with $\mb{s}=\begin{bmatrix}
        s_1 & s_2 & \cdots
    \end{bmatrix}^\mT$ and $\mb{m}=\begin{bmatrix}
        m_1 & m_2 & \cdots
    \end{bmatrix}^\mT$ and applying the Sherman--Morrison--Woodbury formula gives
\begin{gather}
    \begin{split}
        \dfrac{1}{1+\sum\dfrac{m_j}{s_j+i\omega}} & =\left(1+\mb{1}^\mT\diag{\mb{s}+i\omega}^{-1}\mb{m}\right)^{-1} \\&=1-\mb{1}^\mT\left(\diag{\mb{s}+i\omega}+\mb{m}\mb{1}^\mT\right)^{-1}\mb{m}.
    \end{split}
\end{gather}
Assuming the matrix $\diag{\mb{s}}+\mb{m}\mb{1}^\mT$ is neither defective nor degenerate and denoting its eigenvalues as $s'_j$ and the corresponding left and right eigenvectors as $\mb{\phi}^l_j$ and $\mb{\phi}^r_j$, then
\begin{gather}
    \dfrac{1}{1+\sum\dfrac{m_j}{s_j+i\omega}}=1+\sum\dfrac{m'_j}{s_j'+i\omega}\qquad\text{with}\quad{}m'_j=-\dfrac{\left\langle\mb{1},\mb{\phi}^r_j\right\rangle\left\langle\mb{\phi}^l_j,\mb{m}\right\rangle}{\left\langle\mb{\phi}^r_j,\mb{\phi}^l_j\right\rangle}.
\end{gather}

For suitable sets of $s_j$ and $m_j$, it is thus possible to find the corresponding sets $s'_j$ and $m'_j$ to convert between \eqsref{eq:acc_cha} and \eqsref{eq:disp_cha}.
This process naturally extends the definition of $s_j$ and $m_j$ to the complex domain.
Even if the input pairs are real, for example, $s_j,m_j\in\mathbb{R}$, the corresponding $s'_j$ and $m'_j$ can be complex because $s'_j$, which are eigenvalues of the matrix $\diag{\mb{s}}+\mb{m}\mb{1}^\mT$, are not necessarily real.
In the absence of other energy-dissipating actions, the two formulations are thus equivalent.
This equivalence may offer numerical advantages when the mass matrix is constant, as is often the case, and more sparse than the stiffness matrix, particularly under a lumped mass formulation.
Note that this equivalence is only valid under free vibration.
For forced vibration, the external load $\bp$ must be properly transformed to ensure the same equivalence, see later sections.
\subsection{Wrong Asymptotic Behaviour}\label{sec:fix}
The fundamental consistency requirement states that $\mb{K}^*$ (or $\mb{K}'$) must approach the static stiffness $\mb{K}$ at low frequencies, that is, $\lim\limits_{\omega\rightarrow0}\mb{K}^*=\mb{K}$, to ensure the system behaves as expected under static loading conditions.
However, $\mb{K}'$ fails to meet this requirement in \eqsref{eq:u_stiffnesses}.
The formulation must therefore be augmented to fix this issue, for example, by introducing non-trivial terms $\bh_j$ to cancel the contribution of $\sum\dfrac{m_j}{s_j}\mb{K}$ at low frequencies, resulting in
\begin{gather}\label{eq:u_right}
    \left\{\begin{array}{l}
        \mb{M}\ddot{\bu}+\mb{K}\bu+\sum\bbf_j+\sum\bh_j=\mb{0}, \\
        \dot{\bbf}_j=-s_j\bbf_j+m_j\mb{K}\bu,\\
s_j\bh_j=-m_j\mb{K}\bu.
    \end{array}\right.
\end{gather}
The corresponding dynamic stiffness $\mb{K}^*$ is
\begin{gather}
    \mb{K}^*=\left(1-\sum\dfrac{m_j}{s_j}+\sum\dfrac{m_j}{s_j+i\omega}\right)\mb{K}=\left(1-\sum\dfrac{m_j}{s_j}\dfrac{i\omega}{s_j+i\omega}\right)\mb{K}
\end{gather}
with storage stiffness $\mb{K}'$ and loss stiffness $\mb{K}''$
\begin{gather}
    \mb{K}'=\left(1-\sum\dfrac{m_j}{s_j}+\sum\dfrac{s_jm_j}{s_j^2+\omega^2}\right)\mb{K},\qquad
    \mb{K}''=-\sum\dfrac{m_j\omega}{s_j^2+\omega^2}\mb{K}.
\end{gather}
With this modification,
\begin{gather}
    \lim\limits_{\omega\rightarrow0}\mb{K}'=\mb{K},\qquad
    \lim\limits_{\omega\rightarrow\infty}\mb{K}'=\left(1-\sum\dfrac{m_j}{s_j}\right)\mb{K}.
\end{gather}

Due to equivalence, the formulation in \secref{sec:a_dep} also does not meet the consistency requirement.
A similar modification shall be applied.
Only the governing system is presented below, without repeating the lengthy derivations.
\begin{gather}\label{eq:a_right}
    \left\{\begin{array}{l}
        \mb{M}\ddot{\bu}+\mb{K}\bu+\sum\bbf_j+\sum\bh_j=\mb{0}, \\
        \dot{\bbf}_j=-s_j\bbf_j+m_j\mb{M}\ba,\\
s_j\bh_j=-m_j\mb{M}\ba.
    \end{array}\right.
\end{gather}
\subsection{Formulation 3}\label{sec:v_dep_k}
We now consider $\by_j=\mb{K}\bv$, for which the system becomes
\begin{gather}\label{eq:vel_formulation}
    \left\{\begin{array}{l}
        \mb{M}\ddot{\bu}+\mb{K}\bu+\sum\bbf_j=\mb{0}, \\
        \dot{\bbf}_j=-s_j\bbf_j+m_j\mb{K}\dot{\bu}.
    \end{array}\right.
\end{gather}
The characteristic equation can be derived as
\begin{gather}\label{eq:vel_cha}
    \det\left(\left(1+\sum\dfrac{m_ji\omega}{s_j+i\omega}\right)\mb{K}-\omega^2\mb{M}\right)=0.
\end{gather}
Thus,
\begin{gather}
    \mb{K}'=\left(1+\sum\dfrac{m_j\omega^2}{s_j^2+\omega^2}\right)\mb{K},\qquad
    \mb{K}''=\sum\dfrac{s_jm_j\omega}{s_j^2+\omega^2}\mb{K}.
\end{gather}
The limiting cases are
\begin{gather}
    \lim\limits_{\omega\rightarrow0}\mb{K}'=\mb{K},\qquad
    \lim\limits_{\omega\rightarrow\infty}\mb{K}'=\left(1+\sum{}m_j\right)\mb{K}.
\end{gather}
In this formulation, a dissipative system requires $\Re\left(m_j\right)>0$.
A careful comparison reveals that \eqsref{eq:vel_formulation} is equivalent to the model by \citet{Huang2019}, which adopts real parameters $s_j$ and $m_j$.
The loss stiffness $\mb{K}''$ in this case is positive, indicating that the system is dissipative and stable when $\Re\left(m_j\right)>0$.
Unlike the previous formulations, the storage stiffness $\mb{K}'$ approaches the original stiffness $\mb{K}$ at low frequencies and hardens at high frequencies.
It is also equivalent to \eqsref{eq:u_right} after a simple parameter mapping: $s_j$ stays unchanged, and replacing $m_j$ in \eqsref{eq:vel_formulation} by $-m_j/s_j$ yields \eqsref{eq:u_right}.
\subsection{Formulation 4}\label{sec:v_dep_m}
Alternatively, for $\by_j=\mb{M}\dot{\ba}$, the corresponding characteristic equation can be expressed as
\begin{gather}\label{eq:vel_cha2}
    \det\left(\mb{K}-\left(1+\sum\dfrac{m_ji\omega}{s_j+i\omega}\right)\omega^2\mb{M}\right)=0,
\end{gather}
resulting in the dynamic stiffness
\begin{gather}
    \mb{K}^*=\dfrac{1}{1+\sum\dfrac{m_ji\omega}{s_j+i\omega}}\mb{K}.
\end{gather}
As in the previous derivation, partial fraction decomposition can be used to find $s'_j$ and $m'_j$ (explicit expressions are not shown for brevity) such that
\begin{gather}\label{eq:rate_eqv}
    \dfrac{1}{1+\sum\dfrac{m_ji\omega}{s_j+i\omega}}=1+\sum\dfrac{m'_ji\omega}{s'_j+i\omega}.
\end{gather}
%In specific,
%\begin{gather}
%    \begin{split}
%        \dfrac{1}{1+\sum\dfrac{m_ji\omega}{s_j+i\omega}} &=\dfrac{1}{1+\sum{}m_j-\sum\dfrac{m_js_j}{s_j+i\omega}} =\left(1+\sum{}m_j-\mb{s}^\mT\diag{\mb{s}+i\omega}^{-1}\mb{m}\right)^{-1} \\&=\dfrac{1}{\left(1+\sum{}m_j\right)^2}\left(1+\sum{}m_j+\mb{s}^\mT\left(\diag{\mb{s}+i\omega}-\dfrac{\mb{m}\mb{s}^\mT}{1+\sum{}m_j}\right)^{-1}\mb{m}\right).
%    \end{split}
%\end{gather}
%Denoting the eigenvalues of $\diag{\mb{s}}-\dfrac{\mb{m}\mb{s}^\mT}{1+\sum{}m_j}$ as $s'_j$ and the corresponding left and right eigenvectors as $\mb{\phi}^l_j$ and $\mb{\phi}^r_j$,
%\begin{gather}
%m'_js'_j\left(1+\sum{}m_j\right)^2=-\dfrac{\left\langle\mb{s},\mb{\phi}^r_j\right\rangle\left\langle\mb{\phi}^l_j,\mb{m}\right\rangle}{\left\langle\mb{\phi}^r_j,\mb{\phi}^l_j\right\rangle}.
%\end{gather}
Thus, \eqsref{eq:vel_cha2} is an equivalent alternative formulation of \eqsref{eq:vel_cha}.
Because computing and storing $\dot{\ba}=\dddot{\bu}$ is typically \textbf{not} supported in analysis frameworks, this practically restrictive formulation is presented solely to establish the equivalence between the two forms.
If $\dddot{\bu}$ is tracked, the system can leverage the unique characteristics of the mass matrix detailed in previous sections.
\subsection{Forced Vibration}
If the external load $\bp$ is non-trivial, consider the system \eqsref{eq:a_right} with $\bh_j$ merged and an augmented right-hand side,
\begin{gather}
    \left\{\begin{array}{l}
        \left(1-\sum\dfrac{m_j}{s_j}\right)\mb{M}\ddot{\bu}+\mb{K}\bu+\sum\bbf_j=\left(1-\sum\dfrac{m_j}{s_j}\right)\bp+\sum\bq_j, \\
        \dot{\bbf}_j=-s_j\bbf_j+m_j\mb{M}\ba,                                                                                      \\
        \dot{\bq}_j=-s_j\bq_j+m_j\bp.
    \end{array}\right.
\end{gather}
Assuming harmonic solutions for $\bu=\mb{U}\left(\omega\right)\exp\left(i\omega{}t\right)$, $\bbf_j=\mb{F}_j\left(\omega\right)\exp\left(i\omega{}t\right)$, $\bq_j=\mb{Q}_j\left(\omega\right)\exp\left(i\omega{}t\right)$ and $\bp=\mb{P}\left(\omega\right)\exp\left(i\omega{}t\right)$, the equation of motion can be expressed in the frequency domain as
\begin{gather}\label{eq:forced_freq}
    \left(\mb{K}-\left(1-\sum\dfrac{m_j}{s_j}+\sum\dfrac{m_j}{s_j+i\omega}\right)\omega^2\mb{M}\right)\mb{U}=\left(1-\sum\dfrac{m_j}{s_j}+\sum\dfrac{m_j}{s_j+i\omega}\right)\mb{P}.
\end{gather}
It is then straightforward to rewrite \eqsref{eq:forced_freq} as
\begin{multline}
    \left(\dfrac{1}{1-\sum\dfrac{m_j}{s_j}\dfrac{i\omega}{s_j+i\omega}}\mb{K}-\omega^2\mb{M}\right)\mb{U}=\mb{P}\\\longrightarrow\qquad\left(\left(1-\sum\dfrac{m'_j}{s'_j}\dfrac{i\omega}{s'_j+i\omega}\right)\mb{K}-\omega^2\mb{M}\right)\mb{U}=\mb{P}
\end{multline}
using \eqsref{eq:rate_eqv}.
This corresponds to the system \eqsref{eq:u_right} under forced vibration, which is
\begin{gather}
    \left\{\begin{array}{l}
        \mb{M}\ddot{\bu}+\left(1-\sum\dfrac{m_j}{s_j}\right)\mb{K}\bu+\sum\bbf_j=\mb{p}, \\
        \dot{\bbf}_j=-s_j\bbf_j+m_j\mb{K}\bu.
    \end{array}\right.
\end{gather}
\subsection{Remarks}
Because the trivial initial condition $\bbf_j\left(0\right)=\mb{0}$ is assumed, the relation discussed by \citet{Chang2024} can be used to express the proposed formulations equivalently as
\begin{gather}
    \br+\sum{}g_j*\by_j=\mb{p}
\end{gather}
where $g_j*\by_j$ denotes the convolution of the kernel $g_j$ and $\by_j$; that is,
\begin{gather}\label{eq:conv}
    \left(g_j*\by_j\right)\left(t\right)=\int_0^tg_j\left(t-\tau\right)\by_j\left(\tau\right)\md{\tau},
\end{gather}
with $g_j=g_j\left(t\right)=m_j\exp\left(-s_jt\right)$ and $\by_j$ taking one of the forms discussed above.
This type of equivalence provides an alternative perspective to understand various damping formulations.
The damping characteristics, parametrised by $s_j$ and $m_j$, may be interpreted as modifications to the kernel functions $g_j$, enabling the desired damping response to be realised through kernel design.

Among the formulations above, three practical variants are restated in the nonlinear context as follows.
They are equivalent to one another and yield identical energy dissipation characteristics, provided that the parameters $s_j$ and $m_j$ are properly adjusted.
\paragraph{UDD --- Displacement Dependent}
This formulation requires scaling of the tangent stiffness matrix $\mb{K}$.
From an implementation perspective, this operation might be mitigated by applying the reciprocal of the scaling factor to other components of the equation, namely matrices that already require scaling or vectors that are inexpensive to compute.
Its computational cost is lower than that of \eqsref{eq:v_final}, while remaining comparable to that of \eqsref{eq:a_final}, being potentially slightly higher or lower depending on the implementation details and configuration.
\begin{gather}\label{eq:u_final}
\boxed{
    \left\{\begin{array}{l}
        \br_{\ba}+\br_{\bv}+\left(1-\sum\dfrac{m_j}{s_j}\right)\br_{\bu}+\sum\bbf_j=\mb{p}, \\
        \dot{\bbf}_j=-s_j\bbf_j+m_j\br_{\bu}.
    \end{array}\right.
}
\end{gather}
\paragraph{UDV --- Velocity Dependent}
This formulation also requires scaling of the tangent stiffness matrix $\mb{K}$.
In addition, the non-standard quantity $\mb{K}\bv$ must be computed, assembled, and stored as an additional history variable, resulting in the highest computational cost among the three formulations.
\begin{gather}\label{eq:v_final}
\boxed{
    \left\{\begin{array}{l}
        \br_{\ba}+\br_{\bv}+\br_{\bu}+\sum\bbf_j=\mb{p}, \\
        \dot{\bbf}_j=-s_j\bbf_j+m_j\dot{\br}_{\bu}.
    \end{array}\right.
}
\end{gather}
\paragraph{UDA --- Acceleration Dependent}
This formulation requires scaling of the mass matrix $\mb{M}$.
Since $\mb{M}$ is typically sparser than $\mb{K}$ and must be scaled regardless, the additional computational cost is negligible.
This formulation tracks two internal fields $\bbf_j$ and $\bq_j$ instead of one in the other two formulations.
\begin{gather}\label{eq:a_final}
\boxed{
    \left\{\begin{array}{l}
        \left(1-\sum\dfrac{m_j}{s_j}\right)\br_{\ba}+\br_{\bv}+\br_{\bu}+\sum\bbf_j=\left(1-\sum\dfrac{m_j}{s_j}\right)\mb{p}+\sum\bq_j, \\
        \dot{\bbf}_j=-s_j\bbf_j+m_j\br_{\ba},                                                                                      \\
        \dot{\bq}_j=-s_j\bq_j+m_j\bp.
    \end{array}\right.
}
\end{gather}